\documentclass[10pt]{article}

% ==================================================================
% Video Processing (0512-4263, Prof. Shai Avidan, TAU)
% COMPLETE course algorithm reference -- all chapters (1-12) of the
% course html-book. Every algorithm uses the six-part template.
% Every claim is grounded in the provided course material;
% anything not in the source is marked [not specified in source].
% Pure course reference, ordered by chapter.
% ==================================================================

\usepackage[T1]{fontenc}
\usepackage[utf8]{inputenc}
\usepackage[a4paper,margin=2.2cm]{geometry}
\usepackage{amsmath,amssymb,amsfonts}
\usepackage{booktabs}
\usepackage{array}
\usepackage[ruled,vlined,linesnumbered]{algorithm2e}
\usepackage{xcolor}
\usepackage{mdframed}
\usepackage{enumitem}
\usepackage{titlesec}
\usepackage[colorlinks=true,linkcolor=blue!55!black,urlcolor=blue!55!black,
            citecolor=blue!55!black]{hyperref}

\setlength{\parskip}{0.3em}
\setlength{\parindent}{0pt}
\renewcommand{\arraystretch}{1.2}
\setlist{nosep,leftmargin=1.4em}

% six-part run-in headers (compact)
\newcommand{\Purpose}{\par\smallskip\textbf{Purpose.}\ }
\newcommand{\WhenUse}{\par\textbf{When to use.}\ }
\newcommand{\WhenAvoid}{\par\textbf{When to avoid.}\ }
\newcommand{\TheAlgo}{\par\smallskip\textbf{The algorithm.}\par\nopagebreak}
\newcommand{\Knobs}{\par\smallskip\textbf{Parameters and knobs.}\par\nopagebreak}
\newcommand{\ShortExp}{\par\smallskip\textbf{Short explanation.}\ }

% flag box for content absent from the source
\newmdenv[linecolor=red!60!black,backgroundcolor=red!5,linewidth=1pt,
          innertopmargin=5pt,innerbottommargin=5pt]{flagbox}
\newmdenv[linecolor=blue!40!black,backgroundcolor=blue!4,linewidth=0.7pt,
          innertopmargin=5pt,innerbottommargin=5pt]{notebox}

\titleformat{\section}{\normalfont\Large\bfseries\color{blue!30!black}}{\thesection}{0.6em}{}
\titlespacing*{\section}{0pt}{1.4em}{0.6em}
\titleformat{\subsection}{\normalfont\large\bfseries\color{black!85}}{\thesubsection}{0.6em}{}

\newcommand{\Ix}{I_x}\newcommand{\Iy}{I_y}\newcommand{\It}{I_t}
\newcommand{\R}{\mathbb{R}}
\newcommand{\pt}[1]{\marginpar{}}

% smaller font inside parameter tables
\newenvironment{ptable}{\begin{center}\small\begin{tabular}{@{}p{2.5cm} p{5.9cm} p{6.2cm}@{}}\toprule
Parameter & What it controls & Effect of increasing / decreasing \\ \midrule}
{\bottomrule\end{tabular}\end{center}}

\title{\vspace{-1.3cm}\textbf{Video Processing --- Complete Algorithm Reference}\\[2mm]
\large Course 0512-4263 \; (Prof.\ Shai Avidan, Tel-Aviv University)}
\author{}
\date{Compiled from the course html-book, chapters 1--12}

\begin{document}
\maketitle
\vspace{-0.8cm}

\begin{notebox}
\textbf{Scope.} This is a reference to \emph{every} algorithm/method taught in the course
(chapters 1--12), each in a fixed six-part template: \emph{Purpose $\to$ When to use $\to$ When to
avoid $\to$ The algorithm (pseudocode + math) $\to$ Parameters $\to$ Short explanation}. Sections follow
the course chapter order. Every formula is grounded in the source; items not present are marked
\textbf{[not specified in source]}.
\end{notebox}

\tableofcontents

% ##################################################################
\newpage
\section{Motion Representation (Chapter 1)}
% ##################################################################
Every motion a camera perceives is a matrix acting on coordinates. The models form a ladder of
increasing expressiveness (and estimation difficulty): linear $\to$ affine $\to$ homography $\to$ full
3D projection. \emph{Counting rule:} a model with $k$ free parameters needs at least $\lceil k/2\rceil$
point pairs (each pair gives 2 equations); with more pairs the system is over-determined and solved by
Least Squares $x=(A^TA)^{-1}A^Tb$.

% -----------------------------------------------------------------
\subsection{Linear Transformation (2\texttimes2)}
\Purpose Represent rotation, scale, and shear of image coordinates by a $2\times2$ matrix, $p'=Ap$.
\WhenUse Small/far objects whose motion has no translation or perspective; the cheapest model (4 params).
\WhenAvoid Whenever translation is needed --- a pure $2\times2$ matrix always fixes the origin
($A\,(0,0)^T=(0,0)^T$); use affine instead.
\TheAlgo
$\begin{pmatrix}x'\\y'\end{pmatrix}=\begin{pmatrix}a&b\\c&d\end{pmatrix}\begin{pmatrix}x\\y\end{pmatrix}$.
Stack point pairs into $\tilde A\tilde x=b$ with $\tilde x=(a\,b\,c\,d)^T$; each pair gives 2 rows:
\[
\begin{pmatrix} x & y & 0 & 0\\ 0 & 0 & x & y\\ \vdots&&&\vdots\end{pmatrix}
\begin{pmatrix}a\\b\\c\\d\end{pmatrix}=\begin{pmatrix}x'\\y'\\\vdots\end{pmatrix}.
\]
Solve exactly with 2 pairs (invert $4\times4$) or over-determined with LS.
\Knobs
\begin{ptable}
$a,b,c,d$ & Rotation/scale/shear jointly & More free params $=$ more flexible but more noise-sensitive \\
\# point pairs & Determinacy of the fit & $\geq2$ needed; more pairs $\Rightarrow$ LS averages out noise \\
\end{ptable}
\ShortExp ``Multiply by a $2\times2$ matrix.'' Its four numbers encode rotation, scaling and shear
together, but it can never shift the image because the origin is a fixed point. Estimate it from
before/after point pairs (2 equations each); with extra pairs the system is over-determined and Least
Squares gives the best average fit.

% -----------------------------------------------------------------
\subsection{Affine Transformation (2\texttimes3)}
\Purpose Linear transform \emph{plus} translation, using homogeneous coordinates; preserves parallelism.
\WhenUse Medium/closer objects; layer motion (Ch.\ 6 clusters affine motions). 6 params, keeps parallel
lines parallel.
\WhenAvoid When true perspective is present (converging parallels) --- use homography; or when 6 params
overfit a tiny target (use translation/similarity).
\TheAlgo Append a homogeneous ``1'': $\begin{pmatrix}x'\\y'\end{pmatrix}=(A\ t)\begin{pmatrix}x\\y\\1\end{pmatrix}
=\begin{pmatrix}a&b&c\\d&e&f\end{pmatrix}\begin{pmatrix}x\\y\\1\end{pmatrix}$. Each pair gives 2 rows in
6 unknowns $\Rightarrow$ need $\geq3$ pairs; solve $\tilde A\tilde x=b$ by LS.
\Knobs
\begin{ptable}
$A$ (4 params) & Rotation/scale/shear & As linear model \\
$t=(t_x,t_y)$ & Translation & Shifts the origin anywhere (the homogeneous ``1'') \\
\# point pairs & Determinacy & $\geq3$ pairs ($=6$ equations) \\
\end{ptable}
\ShortExp The homogeneous ``1'' gives the matrix a slot to add a constant, so it can translate as if by
multiplication. Affine does rotation $+$ scale $+$ shear $+$ translation and always keeps parallel lines
parallel (a square becomes a parallelogram, never a trapezoid). It needs 3 point pairs.

% -----------------------------------------------------------------
\subsection{Homography (3\texttimes3)}
\Purpose Model a planar perspective mapping (a flat world plane $\to$ image plane), up to scale.
\WhenUse Panorama stitching, document ``flattening'', AR markers, mapping a tilted plane to top-down.
Large/close flat objects.
\WhenAvoid Non-planar / full-3D scenes (needs the projection matrix); estimation is noise-sensitive (8
params).
\TheAlgo $\begin{pmatrix}x'\\y'\\1\end{pmatrix}\cong H\begin{pmatrix}x\\y\\1\end{pmatrix}$, where $\cong$
is equality \emph{up to scale} (divide by the third coordinate: $x'=H_1p/H_3p$, $y'=H_2p/H_3p$). Rearranging
gives a homogeneous system $Ax=0$ (with $A$ of size $2N\times9$); 4 point pairs give $8\times9$, and a
normalisation constraint on the last row fixes the overall scale.
\Knobs
\begin{ptable}
$H$ (9 entries, 8 free) & Perspective warp & Overall scale is free (1 redundancy) $\Rightarrow$ normalise \\
\# point pairs & Determinacy & $\geq4$ pairs (8 equations) \\
Normalisation & Numerical stability & Constrains scale; avoids huge values \\
\end{ptable}
\ShortExp Homography is affine plus true perspective: parallel lines may converge (train tracks meeting).
The $3\times3$ matrix is meaningful only after dividing by the last coordinate --- that division creates
the perspective. Only 8 of its 9 numbers are independent, so 4 point pairs suffice.

% -----------------------------------------------------------------
\subsection{Projection Matrix (pinhole camera)}
\Purpose Map a 3D world point $(X,Y,Z)$ to a 2D pixel $(u,v)$ through the pinhole-camera pipeline.
\WhenUse Relating world geometry to pixels: robotics, AR, self-driving, SLAM, photogrammetry.
\WhenAvoid When the lens distortion matters (undistort first) or intrinsics are uncalibrated.
\TheAlgo Three composed stages (world $\to$ camera $\to$ image $\to$ pixels):
\[
\begin{pmatrix}u\\v\\1\end{pmatrix}\cong
\underbrace{\begin{pmatrix}k_u&0&u_0\\0&k_v&v_0\\0&0&1\end{pmatrix}}_{\text{intrinsics (pixels)}}
\underbrace{\begin{pmatrix}1&0&0&0\\0&1&0&0\\0&0&1/f&0\end{pmatrix}}_{\text{pinhole }(1/Z)}
\underbrace{\begin{pmatrix}R_{3\times3}&t_{3\times1}\\0&1\end{pmatrix}}_{\text{extrinsics }R,t}
\begin{pmatrix}X\\Y\\Z\\1\end{pmatrix}=P\begin{pmatrix}X\\Y\\Z\\1\end{pmatrix}.
\]
Perspective comes from triangle similarity $x=Xf/Z,\ y=Yf/Z$ (far $=$ small). $P$ is $3\times4$ (11 free
params: 6 extrinsic $+$ 5 intrinsic).
\Knobs
\begin{ptable}
$R,t$ (extrinsics) & Camera pose (where it is/points) & Change when the camera moves \\
$f,k_u,k_v,u_0,v_0$ (intrinsics) & Lens $+$ sensor mapping & Calibrated once; $f$ sets zoom, $u_0,v_0$ principal point \\
Depth $Z$ & Projected size ($1/Z$) & Larger $Z$: object projects smaller \\
\end{ptable}
\ShortExp A world point is rotated/translated into camera coordinates ($R,t$), projected onto the image
plane by pinhole similarity (the $1/Z$ that shrinks distant things), then converted from cm to pixels
($k_u,k_v,u_0,v_0$). Multiplying the three matrices gives one $3\times4$ projection matrix $P$ that turns
any $(X,Y,Z,1)$ into a pixel. Extrinsics say where the camera is; intrinsics say how it maps rays to pixels.

% ##################################################################
\newpage
\section{Optical Flow (Chapter 2)}
% ##################################################################
Optical flow asks, for every pixel, where it moved in the next frame: find $(u,v)$ with
$I_1(x{+}u,y{+}v)=I_2(x',y')$. It rests on the \emph{constant brightness constraint} (a moving point keeps
its intensity), which breaks under specularities, shadows, and auto-exposure.

% -----------------------------------------------------------------
\subsection{Lucas--Kanade}\label{sub:lk}
\Purpose Recover local motion $(u,v)$ (or a warp $W(X;P)$) by assuming a whole window moves together and
solving a tiny least-squares system from a first-order Taylor expansion.
\WhenUse Small motions, textured/corner patches, high frame rate; tracking, stabilisation, SLAM --- sparse
flow where speed matters (a $2\times2$ inverse per window).
\WhenAvoid Large jumps; flat/edge-only patches (aperture problem, $A^TA$ singular); specular reflections;
when dense per-pixel flow is needed (use Horn--Schunck).
\TheAlgo Brightness constancy $+$ Taylor gives $\Ix u+\Iy v+\It=0$ per pixel. Stack the window into
$Ax=b$ and minimise $E=\|Ax-b\|^2$:
\[
x=(A^TA)^{-1}A^Tb,\qquad
A^TA=\begin{pmatrix}\sum \Ix^2&\sum \Ix\Iy\\ \sum \Ix\Iy&\sum \Iy^2\end{pmatrix}.
\]
General warp form: $P^*=\arg\min_P\sum_X[I_2(W(X;P))-I_1(X)]^2$, solved incrementally with
$\Delta P^*=-(B^TB)^{-1}B^T\It$, $B=\nabla I_2\,\partial W/\partial P$ (for affine
$\partial W/\partial P=\big(\begin{smallmatrix}x&y&1&0&0&0\\0&0&0&x&y&1\end{smallmatrix}\big)$).

\begin{algorithm}[H]
\SetAlgoLined
\KwIn{$I_1,I_2$; window $W$; warp model $W(X;P)$; tolerance $\epsilon$}
\KwOut{Motion $(u,v)$ / warp params $P$}
Initialise $P$ (identity / zero)\;
\Repeat{$\|\Delta P\|<\epsilon$}{
  Warp $I_2$ by $P$; $\It=I_2(W(X;P))-I_1(X)$\;
  $B=\nabla I_2\,\partial W/\partial P$ over the window\;
  $\Delta P^*=-(B^TB)^{-1}B^T\It$;\quad $P\leftarrow P+\Delta P^*$\;
}
\caption{Incremental Lucas--Kanade}
\end{algorithm}

\Knobs
\begin{ptable}
Window size $n$ & \# equations per solve & Larger: robust but mixes motions; smaller: worse aperture problem \\
Warp model & Translation (2) vs.\ affine (6) & Richer fits larger windows, needs more data \\
Pyramid levels & Max recoverable motion & More levels handle bigger motion, coarse-to-fine \\
\end{ptable}
\ShortExp Instead of searching for matching pixels, LK uses calculus: intensity change $\approx$
gradient $\cdot$ motion. One pixel is one equation but there are two unknowns, so LK assumes a window shares
one motion, giving an over-determined system solved by a $2\times2$ inverse (fast). It fails on flat/edge
patches (rank-deficient) and large motion, which an image pyramid fixes.

% -----------------------------------------------------------------
\subsection{Harris Corner Detector}\label{sub:harris}
\Purpose Find distinctive ``interest points'' (corners) --- exactly the pixels where LK is rank-2 and
trackable.
\WhenUse As the front-end of tracking: detect corners on frame 1, track them with LK (SLAM, 3D
reconstruction).
\WhenAvoid Textureless regions; when dense flow is needed; scale-sensitive (combine with a pyramid).
\TheAlgo The self-difference $E(u,v)=\sum_W[I(x{+}u,y{+}v)-I(x,y)]^2$ linearises to the quadratic form
\[
E(u,v)=(u\ v)\,H\begin{pmatrix}u\\v\end{pmatrix},\quad
H=\sum_{W}\begin{pmatrix}\Ix^2&\Ix\Iy\\ \Iy\Ix&\Iy^2\end{pmatrix}\ (=A^TA),
\]
with eigenvalues $\lambda_{\min}\le\lambda_{\max}$: rank 0 flat, 1 edge, 2 corner.

\begin{algorithm}[H]
\SetAlgoLined
\KwIn{Image $I$; window $g$; sensitivity $\alpha$; threshold $\theta$}
\KwOut{Corner keypoints}
$(\Ix,\Iy)=\mathrm{grad}(I)$;\quad $I_{xx}=\Ix\!\cdot\!\Ix,\ I_{xy}=\Ix\!\cdot\!\Iy,\ I_{yy}=\Iy\!\cdot\!\Iy$\;
$S_{xx}=\mathrm{conv}(I_{xx},g)$, etc.\ ($g=\mathrm{ones}(5,5)$)\;
$R\leftarrow$ one of: $\lambda_{\min}$ (Shi--Tomasi); $\det H/\mathrm{tr}\,H$ (Noble);
$\det H-\alpha(\mathrm{tr}\,H)^2$ (Harris)\;
$R(R<\theta)\leftarrow0$;\quad return max $R$ per tile (non-max suppression)\;
\caption{Harris / structure-tensor detection}
\end{algorithm}
with $\det H=\lambda_{\min}\lambda_{\max}=S_{xx}S_{yy}-S_{xy}^2$, $\mathrm{tr}\,H=S_{xx}+S_{yy}$.

\Knobs
\begin{ptable}
$\alpha$ ($0.04$--$0.06$) & Edge penalty in Harris $R$ & Larger: rejects more edges; smaller: more permissive \\
Threshold $\theta$ & Minimum cornerness & Higher: fewer, stronger corners \\
Window $g$ & Region summed into $H$ & Larger: smoother, rotation-robust \\
\end{ptable}
\ShortExp A pixel on a uniform wall is untrackable. Harris reuses LK's matrix $H=A^TA$ (the structure
tensor) and its eigenvalues to classify pixels: two large $=$ corner, one large $=$ edge, both small $=$
flat. The response $R$ compresses this into one number (Harris's form is cheapest, avoiding
eigen-decomposition); thresholding $+$ non-max suppression keeps strong, spread-out corners.

% -----------------------------------------------------------------
\subsection{Horn--Schunck}\label{sub:hs}
\Purpose Estimate a \emph{dense} flow field for every pixel by combining the brightness-constancy data
term with a global smoothness prior.
\WhenUse Dense flow (warping, interpolation, magnification); coherent motion; small motion.
\WhenAvoid Large motion (use a pyramid); sharp motion boundaries (quadratic smoothness blurs them);
brightness-constancy violations.
\TheAlgo Minimise over the whole image
\[
(u,v)=\arg\min_{u,v}\iint\big[(\Ix u+\Iy v+\It)^2+\alpha(|\nabla u|^2+|\nabla v|^2)\big]\,dx\,dy .
\]
Euler--Lagrange gives $(\Ix u+\Iy v+\It)\Ix-\alpha(u_{xx}+u_{yy})=0$ (and the $v$ analogue). Using the
Laplacian kernel $L=\big(\begin{smallmatrix}0&-1&0\\-1&4&-1\\0&-1&0\end{smallmatrix}\big)$:
\[
\begin{pmatrix}\Ix^2+\alpha L&\Ix\Iy\\ \Ix\Iy&\Iy^2+\alpha L\end{pmatrix}\begin{pmatrix}u\\v\end{pmatrix}
=\begin{pmatrix}\Ix\It\\ \Iy\It\end{pmatrix},
\]
a sparse $2N\times2N$ system ($\approx7$ non-zeros/row), solved iteratively.

\begin{algorithm}[H]
\SetAlgoLined
\KwIn{$I_1,I_2$; weight $\alpha$; iterations $K$}
\KwOut{Dense flow $(u,v)$}
Compute $\Ix,\Iy,\It$; init $u,v\leftarrow0$\;
\For{$k=1$ \KwTo $K$}{
  neighbourhood averages $\bar u,\bar v$\;
  $u\leftarrow\bar u-\Ix\dfrac{\Ix\bar u+\Iy\bar v+\It}{\alpha+\Ix^2+\Iy^2}$;\quad
  $v\leftarrow\bar v-\Iy\dfrac{\Ix\bar u+\Iy\bar v+\It}{\alpha+\Ix^2+\Iy^2}$\;
}
\caption{Horn--Schunck iterative solver}
\end{algorithm}

\Knobs
\begin{ptable}
$\alpha$ (smoothness) & Data vs.\ prior (ML vs.\ MAP) & Larger: smoother, washes out real motion; smaller: choppy/noisy \\
$K$ (iterations) & Solver convergence & More: closer to exact, costlier \\
\end{ptable}
\ShortExp LK works on a window; HS solves the whole image at once and cures the under-determined
``one equation, two unknowns'' problem with a smoothness prior: neighbouring pixels should have similar
flow. $\alpha$ sets the balance (small $=$ trust data / ML, large $=$ trust smooth prior / MAP). The huge
system is very sparse, so it solves fast, but the quadratic penalty blurs true motion boundaries.

% ##################################################################
\newpage
\section{Applications of Motion Estimation (Chapter 3)}
% ##################################################################
Three uses of estimated motion --- redistributing pixel information across time.

% -----------------------------------------------------------------
\subsection{Super-Resolution (ML / MAP)}
\Purpose Fuse several sub-pixel-shifted low-resolution frames into one high-resolution image $X$.
\WhenUse Multiple LR frames of the same scene with \emph{varied} sub-pixel shifts (camera shake helps);
need $\gtrsim r^2$ frames for upscale factor $r$.
\WhenAvoid Identical sub-pixel offsets (no new information); inaccurate motion $F_i$; ML alone (ringing ---
add a prior).
\TheAlgo Forward model per frame: $Y_i=D H F_i X+v_i$ (decimate $D$, blur $H$, warp $F_i$).
\textbf{ML:} $\min_X\sum_i\|DHF_iX-Y_i\|_2^2$, solved by gradient descent
\[
X_{n+1}=X_n-\beta\sum_i F_i^TH^TD^T(DHF_iX_n-Y_i).
\]
\textbf{MAP} (adds a Gibbs/smoothness prior $P(x)\propto e^{-A(x)}$, e.g.\ Laplacian $A(x)=\|Lx\|$):
\[
X_{\text{map}}=\arg\min_X\|DHFX-Y\|_2^2+\lambda A(x),\quad
X_{n+1}=X_n-\beta\!\sum_i F_i^TH^TD^T(DHF_iX_n-Y_i)+\gamma L(X_n).
\]

\begin{algorithm}[H]
\SetAlgoLined
\KwIn{LR frames $\{Y_i\}$; blur $H$, decimation $D$; step $\beta$; (MAP) $\gamma$}
\KwOut{HR image $X$}
Estimate motions $F_i$ ($F_1=$ identity); init $X$ (e.g.\ upsampled $Y_1$)\;
\Repeat{convergence}{
  residual per frame $DHF_iX-Y_i$; back-project by $F_i^TH^TD^T$\;
  $X\leftarrow X-\beta\sum_i F_i^TH^TD^T(DHF_iX-Y_i)$ \ (MAP: $+\gamma L(X)$)\;
}
\caption{Iterative super-resolution (ML/MAP)}
\end{algorithm}

\Knobs
\begin{ptable}
$\beta$ (step) & Gradient-descent rate & Larger: faster but may diverge \\
$\lambda,\gamma$ (prior weight) & Smoothness vs.\ data (MAP) & Larger: smoother, less detail; smaller: sharper, noisier \\
\# frames & Recoverable detail & Need $\gtrsim r^2$ with varied sub-pixel shifts \\
\end{ptable}
\ShortExp Each LR frame samples the HR scene on a slightly different sub-pixel grid, so combining them
recovers detail. The forward model chains warp$\to$blur$\to$downsample; inverting it as least squares is a
huge system, so it is solved by gradient descent using function calls (warp/blur/downsample), not giant
matrices. ML alone rings, so MAP adds a smoothness prior --- the same data-plus-regulariser recipe as HS.

% -----------------------------------------------------------------
\subsection{BRIEF Descriptor}
\Purpose Describe a keypoint with a fast binary ``fingerprint'' for matching.
\WhenUse Fast feature matching (e.g.\ inside hyperlapse); brightness-change resilience; matched by Hamming
distance (XOR $+$ popcount).
\WhenAvoid When rotation invariance is required (BRIEF is not rotation-invariant by default).
\TheAlgo For a set of fixed pixel-pair positions $(p_i,q_i)$ in a window, output bit $=1$ if
$p_i>q_i$ else $0$; concatenate into a binary string. Match two descriptors by Hamming distance.
\Knobs
\begin{ptable}
\# pixel pairs & Descriptor length & More bits: more distinctive, slower \\
Window size & Spatial support & Larger: more context, less localised \\
Pair pattern & Fixed sampling layout & Determines robustness/repeatability \\
\end{ptable}
\ShortExp BRIEF picks fixed pixel pairs around a corner and writes 1 if the first is brighter than the
second, giving a bit string. Because it only compares relative brightness it resists lighting changes, and
matching is a single XOR $+$ popcount --- extremely fast.

% -----------------------------------------------------------------
\subsection{RANSAC}
\Purpose Fit a model robustly when the data contains outliers (wrong matches).
\WhenUse Any fit with outliers: estimating a transform from feature matches, line/plane fitting,
fundamental matrix. Always precede a transform-from-matches LS fit with it.
\WhenAvoid Very high outlier ratios make $K$ explode; needs a sensible inlier threshold $t$.
\TheAlgo Repeat: sample the minimal \#points ($n$: 2 for a line, 3 for affine), fit, count inliers
(residual $<t$), keep the model with the most inliers; refit on all inliers. Number of draws:
\[
K=\frac{\log(1-P)}{\log(1-W^n)},\qquad W=\text{inlier fraction},\ P=\text{target success prob}.
\]

\begin{algorithm}[H]
\SetAlgoLined
\KwIn{Matches; model size $n$; inlier threshold $t$; support $d$; draws $K$}
\KwOut{Robust model}
\For{$K$ times}{
  sample $n$ points; compute model; count inliers ($<t$)\;
  \If{\#inliers $>d$}{keep as candidate}
}
Recompute model using all inliers of the best candidate\;
\caption{RANSAC}
\end{algorithm}

\Knobs
\begin{ptable}
$n$ (sample size) & Points to fit the model & Minimal $n$ maximises all-inlier draw probability \\
$t$ (inlier threshold) & Inlier/outlier cutoff & Larger: more inliers, looser fit; smaller: stricter \\
$K$ (iterations) & Success probability & Set via the formula; higher $P$ or lower $W$ $\Rightarrow$ larger $K$ \\
\end{ptable}
\ShortExp One bad match wrecks a least-squares fit. RANSAC samples the fewest points needed to define a
model, counts how many other points agree, repeats many times, and keeps the model with the most votes,
then refits on the inliers. Using the minimal sample keeps the ``all-inliers'' probability high, and the
formula for $K$ says how many tries guarantee success with probability $P$.

% -----------------------------------------------------------------
\subsection{Hyperlapse (cost $+$ dynamic programming)}
\Purpose Turn a long jittery video into a smooth sped-up one by selecting frames that are smooth
continuations of each other (not uniformly spaced).
\WhenUse Stabilised timelapse from hand-held / wearable footage; when uniform sub-sampling looks jittery.
\WhenAvoid When exact playback timing must be uniform (uniform sub-sampling then wins on timing).
\TheAlgo Match cost between frames $i,j$ (feature points via Harris $+$ BRIEF, transform $T(i,j)$ via
affine $+$ RANSAC): $C_r(i,j)=\frac1n\sum_p\|(x_p,y_p)_j^T-T(i,j)(x_p,y_p)_i^T\|_2$. Full cost adds a
speedup term and a constant-pace term:
\[
C(h,i,j,\mu)=C_m(i,j)+\lambda_s\,C_s(i,j,\mu)+\lambda_a\,C_a(h,i,j),
\]
$C_s=\min(\|(j-i)-\mu\|_2^2,T_s)$, $C_a=\min(\|(j-i)-(i-h)\|_2^2,T_a)$. Minimise the total over paths by a
2-pass dynamic program (build cost matrix, back-trace optimal path).

\begin{algorithm}[H]
\SetAlgoLined
\KwIn{Frames; desired speedup $v$; window $w$; weights $\lambda_s,\lambda_a$}
\KwOut{Frame path $P$}
Init: $D_v(i,j)=C_m(i,j)+\lambda_s C_s(i,j,v)$ for gaps $\le w$\;
Pass I: $D(i,j)=C_m(i,j)+\lambda_sC_s+\min_k[D_v(i{-}k,i)+\lambda_aC_a(i{-}k,i,j)]$; store argmin in $T$\;
Pass II: $P\leftarrow$ Back-Trace$(T)$\;
\caption{Hyperlapse frame selection (DP)}
\end{algorithm}

\Knobs
\begin{ptable}
$\mu$ / $v$ (speedup) & Target frames skipped & Larger: faster video \\
$\lambda_s$ (speedup weight) & Adherence to target rate & Usually $\gg\lambda_a$ (timing matters less than smoothness) \\
$\lambda_a$ (accel.\ weight) & Pace uniformity & Larger: more constant intervals \\
window $w$ & Max jump considered & Larger: more options, costlier DP \\
\end{ptable}
\ShortExp Keeping every $k$-th frame gives violent jitter because samples land at random phases of the
camera's bob. Instead, a cost measures how smoothly frame $j$ follows frame $i$ (via matched features and a
fitted transform), plus penalties for missing the target speedup and for uneven pacing. A two-pass dynamic
program finds the frame-selection path minimising the total cost.

% ##################################################################
\newpage
\section{Online Tracking: Kalman \& Particle Filters (Chapter 4)}
% ##################################################################
Online $=$ frames arrive one at a time. A filter fuses the accumulated belief with each new measurement:
\emph{predict $+$ correct}.

% -----------------------------------------------------------------
\subsection{Bayesian Recursive Tracking (derivation)}\label{sub:bayes}
\Purpose Derive the recursive posterior that underlies both Kalman and particle filters.
\WhenUse Any online tracking where a belief over the state history is updated recursively.
\WhenAvoid When the two assumptions (Markov state; observation depends only on current state) do not hold.
\TheAlgo With state $x_t$, measurements $z_t$: \textbf{Markov} $p(x_t\mid x_{1:t-1})=p(x_t\mid x_{t-1})$
so $p(x_{1:t})=p(x_1)\prod_{j\ge2}p(x_j\mid x_{j-1})$; \textbf{conditional independence}
$p(z_{1:t}\mid x_{1:t})=\prod_\tau p(z_\tau\mid x_\tau)$. Bayes then gives
\[
p(x_{1:t}\mid z_{1:t})\propto
\underbrace{p(z_t\mid x_t)}_{\text{likelihood}}\,
\underbrace{p(x_t\mid x_{t-1})}_{\text{motion model}}\,
\underbrace{p(x_{1:t-1}\mid z_{1:t-1})}_{\text{recursion}},
\]
and marginalising gives filtering
$p(x_t\mid z_{1:t})=C\,p(z_t\mid x_t)\!\int p(x_t\mid x_{t-1})p(x_{t-1}\mid z_{1:t-1})\,dx_{t-1}$.
\Knobs
\begin{ptable}
$p(x_t\mid x_{t-1})$ & Motion / transition model & Constant position/velocity/acceleration \\
$p(z_t\mid x_t)$ & Measurement likelihood & e.g.\ $e^{-\text{SSD}}$ template score \\
prior (recursion) & Carried-forward belief & No need to store the video \\
\end{ptable}
\ShortExp Two assumptions --- the future depends only on the present (Markov), and each observation
depends only on the current state --- collapse Bayes' rule into \emph{likelihood $\times$ motion model
$\times$ previous posterior}. That single recursion is what Kalman and particle filters implement.

% -----------------------------------------------------------------
\subsection{Kalman Filter}\label{sub:kalman}
\Purpose Optimally fuse a linear prediction with a noisy measurement under Gaussian noise.
\WhenUse Roughly linear dynamics, roughly Gaussian noise, unimodal belief; cheap ($O(n^3)$ in state dim).
\WhenAvoid Non-linear dynamics, non-Gaussian noise, or multi-modal beliefs (use a particle filter / EKF).
\TheAlgo Model $x_{k+1}=\Phi x_k+w_k$, $z_k=Hx_k+v_k$, $Q=E[ww^T]$, $R=E[vv^T]$:
\begin{align*}
\text{Predict: } &\hat x'_{k+1}=\Phi\hat x_k, & P'_{k+1}&=\Phi P_k\Phi^T+Q,\\
\text{Gain: } &K_k=P'_kH^T(HP'_kH^T+R)^{-1}, &&\\
\text{Update: } &\hat x_k=\hat x'_k+K_k(z_k-H\hat x'_k), & P_k&=(I-K_kH)P'_k .
\end{align*}

\textbf{Variable meanings:}
\begin{itemize}[noitemsep]
  \item $\hat x_k$ = \textit{state estimate at frame $k$} (after measurement update; this is your best guess right now)
  \item $\hat x'_k$ = \textit{predicted state before measurement} (from the motion model $\Phi$ alone, uncertainty is high)
  \item $z_k$ = \textit{measurement at frame $k$} (e.g., LK optical flow, template match location; noisy data)
  \item $P_k$ = \textit{uncertainty (error covariance) of state estimate}; small = confident, large = uncertain
  \item $P'_k$ = \textit{uncertainty after prediction}; always grows because the motion model drifts
  \item $K_k$ = \textit{Kalman gain} (how much to trust the new measurement vs.\ the prediction; the key dial)
  \item $\Phi$ = \textit{state transition matrix} (e.g., position $\to$ new position, or $[x, v] \to [x+v\Delta t, v]$)
  \item $H$ = \textit{measurement matrix} (extracts what you measured from the state; often identity)
  \item $Q$ = \textit{process noise covariance} (how much does the real world deviate from $\Phi x_k$?)
  \item $R$ = \textit{measurement noise covariance} (how noisy is $z_k$?)
\end{itemize}

$K$ minimises $\mathrm{tr}(P_k)$; in 1-D $K=P'/(P'+R)$ ($R\to0\Rightarrow K\to1$ trust data;
$R\to\infty\Rightarrow K\to0$ trust model).

\begin{algorithm}[H]
\SetAlgoLined
\KwIn{$\Phi,H,Q,R$; init $\hat x_0,P_0$; measurements $z_k$}
\KwOut{Filtered $\hat x_k$}
\For{each frame $k$}{
  Predict: $\hat x'_k=\Phi\hat x_{k-1}$, $P'_k=\Phi P_{k-1}\Phi^T+Q$\;
  Gain: $K_k=P'_kH^T(HP'_kH^T+R)^{-1}$\;
  Measure $z_k$ (LK / template match near $\hat x'_k$)\;
  Update: $\hat x_k=\hat x'_k+K_k(z_k-H\hat x'_k)$, $P_k=(I-K_kH)P'_k$\;
}
\caption{Kalman filter}
\end{algorithm}

\Knobs
\begin{ptable}
$Q$ (process noise) & Trust in motion model & Too small: over-confident, drifts in occlusion; too big: jumpy \\
$R$ (measurement noise) & Trust in sensor & $R$ small $\Rightarrow$ follow data; $R$ large $\Rightarrow$ follow model \\
State vector & What is tracked & Add velocity/accel to coast through occlusions \\
\end{ptable}
\ShortExp Each frame the filter predicts (push the state through $\Phi$; uncertainty grows by $Q$) then
updates (blend the measurement via gain $K$; uncertainty shrinks). $K$ is a trust dial set by $R$ versus the
predicted covariance. Optimal --- but only under linear dynamics and Gaussian noise.

% -----------------------------------------------------------------
\subsection{Particle Filter (Condensation)}\label{sub:particle}
\Purpose Track a non-Gaussian / multi-modal belief as a cloud of $N$ weighted samples.
\WhenUse Non-linear dynamics, non-Gaussian noise, multiple hypotheses (look-alike targets); $100$--$200$
particles.
\WhenAvoid When Kalman's assumptions hold (cheaper); cost is $O(NM)$ per frame; beware sample
impoverishment.
\TheAlgo Samples $\{S_t^{(n)},\pi_t^{(n)},C_t^{(n)}\}$ ($S$ state, $\pi$ weight, $C$ cumulative weight)
implement the recursion of \S\ref{sub:bayes}.

\begin{algorithm}[H]
\SetAlgoLined
\KwIn{$\{S_{t-1}^{(n)},\pi_{t-1}^{(n)},C_{t-1}^{(n)}\}_{n=1}^N$}
\KwOut{New particle set $+$ estimate}
\For{$n=1$ \KwTo $N$}{
  Select: draw $r\sim U[0,1]$; smallest $j$ with $C_{t-1}^{(j)}>r$; $S_t'^{(n)}=S_{t-1}^{(j)}$\;
  Predict: $S_t^{(n)}=A\,S_t'^{(n)}+B\,w_t^{(n)}$ \tcp*{drift $+$ diffuse}
  Measure: $\pi_t^{(n)}=p(z_t\mid x_t=S_t^{(n)})$ \tcp*{e.g.\ $e^{-\text{SSD}}$}
}
Normalise $\sum_n\pi_t^{(n)}=1$; $C_t^{(0)}=0,\ C_t^{(n)}=C_t^{(n-1)}+\pi_t^{(n)}$\;
\Return $\sum_n\pi_t^{(n)}f(S_t^{(n)})$ (or max-weight particle)\;
\caption{Condensation (particle filter)}
\end{algorithm}

\Knobs
\begin{ptable}
$N$ (particles) & Approximation quality & More: better fidelity, costlier; too few: poor coverage \\
Diffusion noise & Exploration & Prevents impoverishment; too much: noisy estimate \\
Output rule & Max vs.\ weighted mean & Max: sharp but jittery; mean: stable but can fall between modes \\
\end{ptable}
\ShortExp A Gaussian cannot represent ``on the table OR under the chair.'' Particles store the belief as
weighted samples of any shape. Each frame: resample (duplicate heavy particles via the cumulative-weight
roulette), predict (motion $+$ noise), re-weight by measurement similarity. Output is a weighted average
(smooth) or the max-weight particle (sharp).

% ##################################################################
\newpage
\section{Dynamic Texture Synthesis (Chapter 5)}
% ##################################################################
\subsection{SVD $+$ ARMA State-Space Texture Synthesis}\label{sub:texture}
\Purpose Learn a low-dimensional linear state-space model of a stochastic texture and drive it with noise
to generate endless novel frames.
\WhenUse ``Dynamic textures'' --- statistically stationary content that never repeats (fire, smoke, water,
grass). State dim $n\sim50$, image dim $m\sim200\times300$.
\WhenAvoid Structured objects (a face, a person walking): the linear low-rank model produces blurry blobs
--- use GAN/VAE.
\TheAlgo Model $x(t{+}1)=Ax(t)+Bv(t)$, $y(t)=Cx(t)+w(t)$. Stack frames $Y_1^\tau=[y(1)\cdots y(\tau)]$; with
$C^TC=I$ take the SVD $Y_1^\tau=U\Sigma V^T\Rightarrow\hat C=U,\ \hat X=\Sigma V^T$; fit dynamics by LS
$\hat A=\arg\min_A\|X_1^\tau-AX_0^{\tau-1}\|$ (AR view: $x_t=C+\sum_i p_ix_{t-i}+\epsilon_t$).

\begin{algorithm}[H]
\SetAlgoLined
\KwIn{Frames $y(1..\tau)$; state dim $n$; noise levels}
\KwOut{Endless synthetic frames}
Vectorise frames into columns $Y$; $U,\Sigma,V^T\leftarrow$ SVD$(Y)$ truncated to $n$\;
$\hat C\leftarrow U$; $\hat X\leftarrow\Sigma V^T$; $\hat A\leftarrow\arg\min_A\|X_1^\tau-AX_0^{\tau-1}\|$\;
pick start state $x$\;
\While{generating}{$x\leftarrow\hat A x+\epsilon$; output $y=\hat C x+w$\;}
\caption{Dynamic-texture learning $+$ synthesis}
\end{algorithm}

\Knobs
\begin{ptable}
$n$ (state dim) & Rank kept from SVD ($\ge95\%$ variance) & Larger: sharper, costlier; smaller: blurrier \\
State noise $v$ & Trajectory exploration & Keeps it non-deterministic; too big: unnatural \\
Image noise $w$ & Natural look & Hides low-rank seams; too big: grainy \\
$C^TC=I$ & Canonical solution & Removes the $T$-ambiguity ($A\to TAT^{-1}$) \\
\end{ptable}
\ShortExp A dynamic texture is a hidden low-dimensional state that drifts by a linear rule and renders to an
image by another linear map --- a Kalman model whose $A,C$ are estimated from data (system identification).
The SVD of stacked frames gives the patterns ($C=U$) and their time-mixing ($x_t=\Sigma V^T$); a least-squares
fit gives $A$. Driving the state forward with noise generates infinitely many plausible frames.

% ##################################################################
\newpage
\section{Layer Representation (Chapter 6)}
% ##################################################################
\subsection{Layer Decomposition (LK $+$ K-means) and Alpha Compositing}
\Purpose Given a composited video and the number of layers $k$, recover each layer, its motion, and its
mask by segmenting the motion into $k$ affine motions.
\WhenUse Videos built from a few coherently-moving layers (background $+$ a few objects); editing, object
removal, background swap. Start with $k=2$--$3$.
\WhenAvoid Non-coherent / highly deformable motion; wrong $k$ (too small lumps motions, too large splits
objects); poor K-means init (empty clusters).
\TheAlgo Alpha-compositing synthesis model $I_k=I_{k-1}(1-\alpha_k)+E_k\alpha_k$ (layers $E_k$, masks
$\alpha_k$). Analysis is an EM-style alternation:

\begin{algorithm}[H]
\SetAlgoLined
\KwIn{Frame pair $I_1,I_2$; \# layers $k$}
\KwOut{Per-pixel layer assignment, per-layer affine motion}
\Repeat{convergence}{
  \textbf{LK:} fit a 6-D affine motion $V\in\R^6$ on random windows\;
  \textbf{K-means:} cluster the window motions in $\R^6$ into $k$ groups (layers)\;
  Refit LK per layer on its pixels\;
}
\caption{Layer decomposition (LK $\to$ K-means loop)}
\end{algorithm}

\Knobs
\begin{ptable}
$k$ (\# layers) & Number of motion groups & Too small: lumps motions; too large: splits objects (MDL picks $k$) \\
Motion model & Affine (6) vs.\ homography (8) & Affine usually enough; homography for planar perspective \\
\# random windows & Sampling of motions & More: accurate, slower; random sub-sampling is standard \\
\end{ptable}
\ShortExp Movies are layered sandwiches (background, characters, effects) combined by alpha compositing
$I=E\alpha+I_{\text{prev}}(1-\alpha)$. To invert this, assume each layer moves by one affine transform (6
numbers): fit LK motions on windows, cluster them in $\R^6$ with K-means (each cluster $=$ a layer), refit
per layer, and iterate --- the classic ``assume A to estimate B, swap, repeat''.

% -----------------------------------------------------------------
\subsection{K-means Clustering}\label{sub:kmeans}
\Purpose Partition points $x\in\R^d$ into $k$ groups minimising total squared-error distortion.
\WhenUse Grouping vectors by similarity (motion vectors into layers; vector quantisation) with known $k$.
\WhenAvoid Unknown $k$ (needs MDL/trial); non-convex/unequal clusters; poor init (local minima).
\TheAlgo Minimise $\text{Distortion}=\sum_i\|x_i-c_{\text{Encode}(x_i)}\|^2$; setting
$\partial/\partial c_j=0$ gives $c_j=\frac{1}{|s_j|}\sum_{i\in s_j}x_i$.

\begin{algorithm}[H]
\SetAlgoLined
\KwIn{Points $\{x_i\}_{i=1}^n$; \# clusters $k$}
\KwOut{Prototypes $c_1,\dots,c_k$}
Initialise $c_1,\dots,c_k$ (random / $k$-means++)\;
\Repeat{assignments stop changing}{
  Assign each $x_i$ to nearest $c_j$\;
  Update $c_j\leftarrow\frac{1}{|s_j|}\sum_{i\in s_j}x_i$\;
}
\caption{K-means (Lloyd)}
\end{algorithm}

\textbf{Termination (guaranteed).} Both steps never increase the distortion (nearest-centre assignment
minimises it for fixed centres; the mean minimises it for fixed assignments), so distortion is monotonically
non-increasing. Since there are only finitely many partitions ($\le k^n$) and the cost cannot cycle upward,
the algorithm stops after finitely many steps --- at a local (not necessarily global) minimum.

\Knobs
\begin{ptable}
$k$ (\# clusters) & Granularity & Larger $k$: lower distortion ($0$ at $k=n$); MDL picks optimal $k$ \\
Initialisation & Which local minimum & $k$-means++ avoids empty clusters \\
Restarts & Robustness & Keep the lowest-distortion run \\
\end{ptable}
\ShortExp Drop $k$ centres, assign each point to its nearest, move each centre to its members' mean,
repeat. Each step is a distortion minimiser, so the cost only ever drops; with finitely many partitions and
no upward cycling, K-means is guaranteed to terminate --- possibly at a local minimum, which is why
initialisation and restarts matter.

% ##################################################################
\newpage
\section{Video Magnification (Chapter 7)}
% ##################################################################
\subsection{Eulerian Video Magnification}
\Purpose Amplify motions too small to see (a pulse, a swaying building) without estimating motion.
\WhenUse Sub-pixel, low-amplitude, dense changes; a steady (tripod) camera; known target frequency band.
\WhenAvoid Motion $>\!\sim1$ pixel (Taylor breaks $\Rightarrow$ ringing --- use pyramid magnification);
shaky camera (shake gets amplified too).
\TheAlgo For $I(x,t)=f(x+\delta(t))$, first-order Taylor gives
$I(x,t)\approx f(x)+\underbrace{\delta(t)\,\partial f/\partial x}_{B(x,t)}$. Adding $\alpha B$ back yields
\[
\tilde I(x,t)=I(x,t)+\alpha B(x,t)\approx f\big(x+(1+\alpha)\delta(t)\big),
\]
so motion is amplified by $\alpha$ --- \emph{without} estimating $\delta$. In practice $B\approx I_2-I_1$
(ideally band-passed in time first).

\begin{algorithm}[H]
\SetAlgoLined
\KwIn{Video frames; amplification $\alpha$; (optional) temporal band-pass}
\KwOut{Motion-magnified video}
\ForEach{location (independently)}{
  temporal signal $\to$ band-pass to the target frequency $\to$ get $B\approx I_2-I_1$\;
  output $\tilde I=I+\alpha B$\;
}
\caption{Eulerian magnification}
\end{algorithm}

\Knobs
\begin{ptable}
$\alpha$ (amplification) & Motion gain (typ.\ 5--50) & Higher: reveals more but more ringing artifacts \\
Band-pass band & Frequencies amplified & Match target (heart $\approx1$\,Hz); cuts noise \\
Pyramid & Handles $>\!1$px motion & Amplify each scale separately to avoid ringing \\
\end{ptable}
\ShortExp Lagrangian methods follow each pixel (need motion estimation); Eulerian methods sit at a fixed
location and watch its \emph{intensity profile} --- the value of that one pixel over time --- change.
A tiny spatial shift $\delta$ changes the pixel value by
$\delta\times$ slope (Taylor); multiply that temporal change by $\alpha$ and add it back and the pixel
``moves'' as if by $(1+\alpha)\delta$ --- amplification with no motion estimation. Band-pass first so only
the target motion (not lighting drift/noise) is amplified.

% ##################################################################
\newpage
\section{Background Subtraction (Chapter 8)}
% ##################################################################

\textbf{Why GMM? (motivation from the intensity profile).} Fix one pixel and plot its value across every
frame of the video --- its \emph{intensity profile}. Ideally it sits at the background colour, jumps when a
moving object passes over it, then returns. In practice a pixel often has \emph{several} typical background
values (e.g.\ a leaf blowing in front of sky), so its profile oscillates between a few clusters even with no
foreground object present. The naive fix is to collect many frames, build a histogram of that pixel's
values, and cluster it (the big clusters = background modes, e.g.\ ``leaf'' and ``sky''; a small outlier
cluster = a passing object). This works, but has three problems: (1) \emph{delay} --- you must wait to
collect enough frames before clustering; (2) \emph{cost} --- re-clustering per pixel is expensive; (3)
\emph{memory} --- you must store the frame history. The \textbf{GMM} below fixes all three by replacing
batch clustering with an \emph{online, incremental} update: each new pixel value nudges an existing Gaussian
or spawns a new one, with no stored history.

\subsection{Per-pixel Gaussian Mixture Model (GMM)}
\Purpose Classify each pixel as background or a moving object online, robust to repetitive background
motion (rustling leaves).
\WhenUse Static-camera surveillance with multimodal backgrounds; low memory ($O(K)$ state per pixel).
\WhenAvoid Sudden lighting changes, camera shake (false positives until it adapts); slow objects get
absorbed; shadows flagged as foreground.
\TheAlgo Each pixel's value is modelled by $K$ Gaussians
$f_x(x\mid\Phi)=\sum_k\omega_k f_{x\mid k}(x\mid k,\theta_k)$, $\theta_k=(\mu_k,\sigma_k)$,
$\sum_k\omega_k=1$. Decision by MAP $\hat k=\arg\max_k\omega_k f_{x\mid k}$. Sort Gaussians by
$\omega_k/\sigma_k$ (frequent $+$ tight $=$ background) and take the top ones with
$B=\arg\min_b(\sum_{k\le b}\omega_k/\sigma_k>T)$ as background. Online update per new value $x_t$:
\[
\omega_{k,t}=(1-\alpha)\omega_{k,t-1}+\alpha M_{k,t},\quad
\mu_t=(1-\rho)\mu_{t-1}+\rho x_t,\quad
\sigma_t^2=(1-\rho)\sigma_{t-1}^2+\rho(x_t-\mu_t)^2,
\]
where $M_{k,t}=1$ if $x_t$ matches Gaussian $k$ (within $2.5\sigma$), else $0$, and
$\rho=\alpha\,p(k\mid x_t)/\hat\omega_{k,t}$.

\begin{algorithm}[H]
\SetAlgoLined
\KwIn{New pixel value $x_t$; $K$ Gaussians; rates $\alpha,\rho$; threshold $T$}
\KwOut{Background / foreground label}
\eIf{$x_t$ within $2.5\sigma$ of some Gaussian $k$}{
  nudge $\mu_k,\sigma_k$ toward $x_t$; increase $\omega_k$\;
}{
  replace lowest-weight Gaussian with one centred at $x_t$ (low $\omega$, high $\sigma$)\;
}
Renormalise weights; sort by $\omega_k/\sigma_k$; background $=$ top until $\sum>T$\;
Label $x_t$ foreground iff its Gaussian is not background\;
\caption{Online GMM background subtraction}
\end{algorithm}

\Knobs
\begin{ptable}
$K$ (Gaussians) & Modes modelled (3--5) & Fewer: miss multimodal BG; more: fit noise \\
$\alpha$ (weight rate) & How fast new patterns become BG & $\approx0.01$; too fast: parked car $\to$ BG; too slow: no adaptation \\
$T$ (BG threshold) & BG fraction (0.7--0.9) & Higher: more count as BG, fewer false alarms \\
Match ($2.5\sigma$) & New-object sensitivity & Tight: detect changes; loose: absorb noise \\
\end{ptable}
\ShortExp Each pixel keeps its own mixture of a few Gaussians (mean colour, width, popularity). A new value
that matches an existing Gaussian nudges it; otherwise it replaces the weakest one (likely a moving object).
Background Gaussians are the frequent \emph{and} tight ones (high $\omega/\sigma$) whose weights sum past
$T$; everything else is foreground. It is fully online and memory-efficient.

% ##################################################################
\newpage
\section{Video Colorization \& Geodesic Distance (Chapter 9)}
% ##################################################################
Propagate a few colour scribbles to every pixel: fast within similar-luminance regions, slow across edges.
The vehicle is \emph{geodesic distance}, computed by Dijkstra or fast sweeping.

% -----------------------------------------------------------------
\subsection{Geodesic Distance Colorization}\label{sub:colorize}
\Purpose Fill the two chrominance channels of a B\&W video from sparse colour scribbles, respecting
luminance edges.
\WhenUse Region-wise-smooth colour (cartoon-like); colour edges coincide with luminance edges; video needs
far fewer scribbles (time axis carries colour).
\WhenAvoid Saturated luminance with varying colour (fire); smooth colour gradients (sunset); specular
highlights.
\TheAlgo Path length weighted by contour-crossing
$L(p)=\int_p|\langle\dot p(s),\nabla Y\rangle|\,dl$, distance $d_L(x,y)=\min_p L(p)$, distance map
$D_k(x)=d_L(x,\Omega_k)$ to scribble $k$. Blend by a distance-decaying convex combination:
\[
c(x)=\frac{\sum_k W(D_k(x))\,c_k}{\sum_k W(D_k(x))},\qquad W(D)=\frac{1}{D^{\,b}},\ b\in[1,6].
\]

\begin{algorithm}[H]
\SetAlgoLined
\KwIn{Luminance $Y$; scribbles $\{\Omega_k,c_k\}$; exponent $b$}
\KwOut{Chrominance $c(x)$ for all pixels}
\ForEach{scribble $k$}{$D_k\leftarrow$ geodesic distance from $\Omega_k$ (Dijkstra / fast sweeping)\;}
\ForEach{pixel $x$}{$c(x)\leftarrow\sum_k D_k(x)^{-b}c_k\,/\,\sum_k D_k(x)^{-b}$\;}
\caption{Geodesic colorization}
\end{algorithm}

\Knobs
\begin{ptable}
$b$ in $1/D^b$ (1--6) & Transition sharpness & Larger: sharper boundaries; smaller: smoother blends \\
Edge weight $W_e=|Y(p)-Y(q)|$ & How edges act as walls & Gradient magnitude blocks colour bleed \\
Scribble placement & Region coverage & Inside regions, not on edges; more in large regions \\
Video neighbourhood & 4 (image) vs.\ 6 (+time) & Time edges let colour flow between frames \\
\end{ptable}
\ShortExp Redefine distance to follow content: staying inside one region costs almost nothing, crossing an
edge costs a lot. For each scribble, compute this geodesic distance to every pixel; each pixel's colour is a
weighted vote of all scribbles with weight $1/D^b$ (nearby scribbles dominate). In video, adding a temporal
edge (6-neighbourhood) lets colour flow through time, so far fewer scribbles are needed.

% -----------------------------------------------------------------
\subsection{Dijkstra's Algorithm}\label{sub:dijkstra}
\Purpose Compute the shortest (geodesic) distance from a source set to every node of a weighted graph.
\WhenUse Exact single-source distance maps on a CPU with a priority queue; $O(N\log N)$.
\WhenAvoid Massively parallel / GPU settings (the priority queue is hard to parallelise --- use fast
sweeping).
\TheAlgo Edge weight $W_e=|Y(x,y)-Y(x',y')|$; path length $L_p=\sum_{e\in P}W_e$.

\begin{algorithm}[H]
\SetAlgoLined
\KwIn{Grid $(\Omega,E)$, weights $W$; source set $\Omega_s$}
\KwOut{$D(x)=d_L(x,\Omega_s)$}
$D|_{\Omega_s}=0$, $D|_{\text{rest}}=\infty$, $Q=\Omega$\;
\While{$Q\neq\varnothing$}{
  $x=\arg\min_{x\in Q}D(x)$ \tcp*{priority queue $O(\log N)$}
  \ForEach{$y\in N_1(x)\cap Q$}{$D(y)\leftarrow\min\{D(y),D(x)+W(x,y)\}$\;}
  $Q\leftarrow Q\setminus\{x\}$\;
}
\caption{Dijkstra (geodesic distance)}
\end{algorithm}
Correctness rests on the \emph{Bellman principle}: any sub-path of an optimal path is optimal.

\Knobs
\begin{ptable}
Edge weights & Cost of crossing gradients & Gradient-magnitude weighting makes edges walls \\
Neighbourhood $N_1$ & 4 (image) / 6 (video) & Larger: more path options, costlier \\
\end{ptable}
\ShortExp Water spreads from the source pixels but flows slowly through high-gradient regions; the pixel
that is ``wet first'' takes that distance. Pop the closest unvisited pixel, relax its neighbours, repeat ---
$O(N\log N)$ with a priority queue. The Bellman principle (sub-paths of optimal paths are optimal) is what
makes the greedy relaxation correct.

% -----------------------------------------------------------------
\subsection{Fast Sweeping (Danielsson)}
\Purpose Compute the geodesic distance map with parallel-friendly raster sweeps instead of a priority
queue.
\WhenUse GPU / cache-friendly settings; often faster than Dijkstra in practice.
\WhenAvoid Highly spiralling geodesics need many sweeps; distances are not Euclidean so $>4$ sweeps may be
required.
\TheAlgo Update each pixel from already-visited neighbours in raster order; 4 sweeps cover the 4 directions;
repeat sweeps until the distance map stops changing (spirals need $\ge4N$ sweeps for $N$ turns). Parallelise
by updating a whole column at a time.
\Knobs
\begin{ptable}
\# sweeps & Directions / turns covered & 4 for simple paths; more for spirals; iterate to convergence \\
Neighbourhood layout & Parallelisation & Column-wise update enables parallel sweeps \\
\end{ptable}
\ShortExp Dijkstra's priority queue is hard to parallelise. Fast sweeping instead runs simple raster scans,
updating each pixel from its already-visited neighbours; four sweeps handle the four directions and more are
added until nothing changes. It is cache- and GPU-friendly and usually faster than Dijkstra.

% ##################################################################
\newpage
\section{Video Matting (Chapter 10)}
% ##################################################################
Extract a foreground and composite it onto a new background without colour bleed:
trimap $\to$ opacity map $\to$ matting. Reuses Ch.\ 9's geodesic machinery with a foreground-probability
cost field.

% -----------------------------------------------------------------
\subsection{Trimap Estimation (KDE $+$ geodesic)}
\Purpose Label each pixel confidently-foreground (1), confidently-background (0), or unsure ($\phi$).
\WhenUse Given FG/BG scribbles; when colour alone is ambiguous (overlapping FG/BG colours) so geodesic
proximity breaks ties.
\WhenAvoid Poorly-placed scribbles (need one per distinct colour region); wrong band radius $\rho$.
\TheAlgo Estimate colour distributions by KDE (Parzen window)
$f(c\mid F)\propto\sum_{x\in\Omega_F}K(c,c(x))$, $K\propto e^{-\|c-c'\|^2/2\sigma^2}$; Bayes gives
$P(F\mid c)=f(c\mid F)/(f(c\mid F)+f(c\mid B))$ (uniform prior). Compute geodesic distances $D_F,D_B$ using
$W(x)=\nabla_x P(F\mid c(x))$ as the field; classify $V_F=\{D_F<D_B\}$; the unsure band is a radius-$\rho$
ball around the boundary $\Delta=\{D_F=D_B\}$:
\[
\alpha_T(x)=1\ \text{on}\ V_F\setminus B_\rho(\Delta),\quad
\alpha_T(x)=0\ \text{on}\ V_B\setminus B_\rho(\Delta),\quad
\alpha_T(x)=\phi\ \text{on}\ B_\rho(\Delta).
\]
\Knobs
\begin{ptable}
$\rho$ (band radius) & Width of unsure band (5--15 px) & Too small: misses soft hair; too big: too much work for stage 2 \\
$\sigma$ (KDE kernel) & Colour-model smoothness & Larger: smoother distribution \\
Scribbles & Colour-region coverage & One per distinct FG/BG colour region \\
\end{ptable}
\ShortExp Place a soft Gaussian ``bell'' on each scribble pixel's colour (KDE) to get FG/BG colour
distributions, convert to $P(F\mid c)$ by Bayes, then compute geodesic distances to FG and BG scribbles.
Classify by whichever is closer, and mark a narrow band around the boundary as ``unsure''. Geodesic
proximity resolves cases where FG and BG share colours.

% -----------------------------------------------------------------
\subsection{Opacity Map (trimap refinement)}
\Purpose Assign each unsure pixel a continuous $\alpha\in[0,1]$.
\WhenUse Inside the narrow uncertain band, where local colour statistics differ from the global model.
\WhenAvoid Outside the band (already confident 0/1).
\TheAlgo Re-estimate FG/BG distributions using only boundary pixels (nearly 1-D along the boundary curve
length $\tau$), then combine local probability with geodesic distance:
\[
\alpha(x)=\frac{w_F(x)}{w_F(x)+w_B(x)},\qquad w_F(x)=\tilde D_F(x)^{-r}\,\tilde P_F(x),\ r\in(0,2].
\]
\Knobs
\begin{ptable}
$r\in(0,2]$ & Distance weighting & Larger: sharper reliance on proximity \\
Local KDE window & Boundary colour model & Re-estimated from boundary pixels (richer local data) \\
\end{ptable}
\ShortExp Near the boundary there are many more FG/BG samples than the sparse scribbles, so re-estimate the
colour models locally. Combining the local FG probability with the geodesic distances gives each unsure pixel
a soft opacity $\alpha$ between 0 and 1.

% -----------------------------------------------------------------
\subsection{Matting (pure-foreground recovery)}
\Purpose Composite onto a new background without dragging the old background's colour into fuzzy pixels.
\WhenUse The uncertain boundary band (hair, fuzzy edges) where a pixel mixes FG and BG colour.
\WhenAvoid Confident FG/BG regions (a plain $\alpha$ blend is fine there).
\TheAlgo A boundary pixel is $c(x)=\alpha f(x)+(1-\alpha)b(x)$ (one equation, two unknowns). In a small
window assume $f,b$ match confident FG/BG neighbours; pick the pair that best explains $c(x)$:
\[
(x_F^*,x_B^*)=\arg\min_{x_F,x_B}\|\alpha(x)c(x_F)+(1-\alpha(x))c(x_B)-c(x)\|^2,\quad
J(x)=\alpha\,c(x_F^*)+(1-\alpha)B(x).
\]
\Knobs
\begin{ptable}
Window size & Neighbour search for $f,b$ & Larger: more candidates, costlier; small windows usually suffice \\
New background $B$ & Composite target & Provided by the user \\
\end{ptable}
\ShortExp A naive blend $J=\alpha c+(1-\alpha)B$ drags the old background's colour onto the new one at
fuzzy pixels. Instead recover the \emph{pure} foreground colour by finding, in a small window, the
confident FG/BG pixel pair that best reproduces the mixed pixel, and composite that pure colour ---
avoiding colour bleed.

% -----------------------------------------------------------------
\subsection{Fast KDE via Low-Rank Kernel}
\Purpose Evaluate $f(x)=\sum_{y}k(x,y)w(y)$ (KDE) in linear rather than quadratic time.
\WhenUse Large images where naive KDE is $O(|\Omega|\,|\Omega'|)$; kernels with fast-decaying singular
values (e.g.\ Gaussian).
\WhenAvoid Genuinely full-rank kernels (no speedup).
\TheAlgo If $K=\sum_{i=1}^p\phi_i\psi_i^T$ (rank $p\ll m,n$), then with $a_i=\sum_y\psi_i(y)w(y)$,
$f(x)=\sum_i a_i\phi_i(x)$, total $O((m+n)p)$ instead of $O(mn)$ (the Fast Multipole Method factorises such
kernels).
\Knobs
\begin{ptable}
$p$ (effective rank) & Approximation vs.\ speed & Smaller $p$: faster, less accurate \\
Kernel choice & Rank / decay & Pick low-rank (fast-decaying singular values) \\
\end{ptable}
\ShortExp Naive KDE sums over every scribble pixel at every output pixel --- quadratic. If the kernel is
(approximately) low rank, the double sum factors into two cheap passes, giving $O((m{+}n)p)$. Gaussians are
formally infinite-rank but their singular values decay, so a small effective rank works.

% ##################################################################
\newpage
\section{Re-targeting \& PatchMatch (Chapter 11)}
% ##################################################################
Resize/edit an image while preserving a natural look, formalised by two losses (completeness $+$ coherence)
and solved fast by PatchMatch.

% -----------------------------------------------------------------
\subsection{Bidirectional-Similarity Re-targeting}
\Purpose Produce a target of new size/aspect that keeps all important source content (completeness) and
invents nothing (coherence).
\WhenUse Content-aware resize, inpainting, reshuffling, object removal/preservation; images and video (3-D
patches).
\WhenAvoid Bad initialisation (start from black $\Rightarrow$ garbage --- use a hierarchical pyramid);
straight lines/objects may bend (add constraints).
\TheAlgo With patch distance $d^2(P,Q)=\|P-Q\|_2^2$:
\[
D_{\text{com}}(S,T)=\tfrac1{N_S}\!\sum_{P\in S}\min_{Q\in T}d^2(Q,P),\quad
D_{\text{coh}}(S,T)=\tfrac1{N_T}\!\sum_{Q\in T}\min_{P\in S}d^2(P,Q),
\]
$D=\alpha D_{\text{com}}+(1-\alpha)D_{\text{coh}}$. Per-pixel optimum is a weighted average of voting source
colours (set derivative to 0). Weighted constraints multiply patches by $W_P$ (high $=$ preserve, $\approx0$
$=$ remove), keeping $D$ convex.

\begin{algorithm}[H]
\SetAlgoLined
\KwIn{Source $S$; target size $\Omega'$; balance $\alpha$}
\KwOut{Target $T$}
Initialise $T^1$ by hierarchical resize (gradual pyramid), \emph{not} black\;
\Repeat{convergence}{
  reset $T^{k+1}\leftarrow0$, weight map $C\leftarrow0$\;
  \textbf{Coherence:} each $Q\in T^k$ finds NN $P\in S$; add $\frac{1-\alpha}{N_T}P$ onto $Q$; $C{+}{=}\frac{1-\alpha}{N_T}$\;
  \textbf{Completeness:} each $P'\in S$ finds NN $Q'\in T$; add $\frac{\alpha}{N_S}P'$ onto $Q'$; $C{+}{=}\frac{\alpha}{N_S}$\;
  $T^{k+1}\leftarrow T^{k+1}/C$ (element-wise)\;
}
\caption{Re-targeting by bidirectional similarity}
\end{algorithm}

\Knobs
\begin{ptable}
$\alpha$ (0--1) & Completeness vs.\ coherence & Higher: keep more content; lower: smoother but riskier (start 0.5) \\
Patch size (5--7) & Context vs.\ sharpness & Larger: more context, blurrier \\
Constraint weights $W_P$ & Preserve / remove & High: protect a face/line; $\approx0$: erase an object \\
Pyramid step ($\sim1.1\times$) & Init quality & Gradual resize avoids bad local minima \\
\end{ptable}
\ShortExp Two symmetric principles: every important source patch should appear in the target
(completeness), and every target patch should exist in the source (coherence). Their weighted sum is
minimised by setting each output pixel to a weighted average of the source patches voting for it. Good
initialisation (a gradual resize pyramid) is essential, and the nearest-neighbour search is the bottleneck
that PatchMatch accelerates.

% -----------------------------------------------------------------
\subsection{PatchMatch}\label{sub:patchmatch}
\Purpose Compute an approximate nearest-neighbour offset field in near real time.
\WhenUse Whenever brute-force $O(N^2)$ patch matching is too slow and matches are spatially coherent (most
natural images).
\WhenAvoid Non-coherent correspondences; when an exact global NN is required.
\TheAlgo Find $f(x)=\arg\min_y\|P_x-Q_{x+f}\|_2^2$ using \emph{propagation} (exploit coherence) $+$
\emph{random search} (escape local minima).

\begin{algorithm}[H]
\SetAlgoLined
\KwIn{Source $S$, target $T$; window $W$; iterations $K$}
\KwOut{NN offset field $f$}
Init $f\sim U[-W/2,W/2]^2$; $d(x)=\|P_x-Q_{x+f(x)}\|_2^2$\;
\For{$k=1$ \KwTo $K$ (alternating raster order)}{
  \ForEach{$x$}{
    Propagation: $f(x)\leftarrow\arg\min_{f\in\{f(x),f(x-e_1),f(x-e_2)\}}\|P_x-Q_{x+f}\|$\;
    Random search: test $f(x)+r$ with shrinking radius; keep if it lowers $d(x)$\;
  }
}
\caption{PatchMatch (1-NN)}
\end{algorithm}

\Knobs
\begin{ptable}
Patch size (5--7) & Context vs.\ sharpness & Larger: more context, blurrier \\
Iterations $K$ (4--6) & NN-field convergence & More: better matches, costlier \\
Window $W$ & Random-search radius & Shrinks each inner step \\
\end{ptable}
\ShortExp If patch $(100,100)$ matches $(200,50)$, its neighbour almost certainly matches with the same
offset --- so neighbours borrow each other's offset instead of searching. PatchMatch alternates propagation
(spread good offsets to causal neighbours, flipping scan direction each pass) with shrinking-radius random
search. Together they converge in $\sim4$--$6$ passes, hundreds of times faster than brute force.

% -----------------------------------------------------------------
\subsection{Applications: Inpainting, Reshuffling, Object Removal}
\Purpose Fill missing pixels, move objects, or delete content --- all as the same re-targeting optimisation
with constraints, solved by PatchMatch.
\WhenUse Scratch/hole filling; rearranging objects; content-aware removal; mixing content from several
sources.
\WhenAvoid Large semantically-structured holes (patch copying can't hallucinate unseen structure).
\TheAlgo \emph{Inpainting:} minimise $D(S,T)$ over the missing region $\Omega_m$ with
$T|_{\Omega\setminus\Omega_m}=S|_{\Omega\setminus\Omega_m}$. \emph{Reshuffling:} fix pasted-object pixels,
fill the rest. \emph{Removal:} set the object's patch weights $\approx0$ so it cannot influence $T$;
\emph{preservation:} weight protected patches (and their NN) high.
\Knobs
\begin{ptable}
Mask $\Omega_m$ & Region to synthesise & Generous mask ($+$few px) avoids edge halos \\
Weights $W_P$ & Preserve vs.\ remove & High: keep; $\approx0$: remove \\
Source set & Where fill comes from & Multiple sources allowed \\
\end{ptable}
\ShortExp All these tasks reduce to ``define which pixels are fixed and which are free, then minimise
completeness $+$ coherence with PatchMatch''. Weighting patches steers the result --- high weight preserves
a face or line, near-zero weight makes an object vanish as neighbours fill its space.

% ##################################################################
\newpage
\section{Intrinsic Images from Sequences (Chapter 12)}
% ##################################################################
\subsection{Weiss Intrinsic Images (median of derivative filters)}
\Purpose From $T$ images of a fixed scene under changing light, recover one reflectance (paint) image and
$T$ illumination (light) images.
\WhenUse A stationary camera/scene with only lighting varying (a day-long webcam); a handful of frames
already suffices.
\WhenAvoid A single image (ill-posed: $T$ equations, $T{+}1$ unknowns); non-sparse lighting; strongly
non-linear camera response (assumed linear).
\TheAlgo In the log domain $i(x,y,t)=r(x,y)+l(x,y,t)$. Assume derivative-filter outputs of the lighting are
sparse (Laplacian $P(x)\propto e^{-\alpha|x|}$). With filters $\{f_n\}$, $o_n=i\star f_n$, the ML estimate of
the filtered reflectance is the \emph{temporal median}
\[
\hat r_n(x,y)=\operatorname*{median}_t\,o_n(x,y,t)
\]
(because maximising the Laplacian likelihood minimises the $\ell_1$ error, whose minimiser is the median).
Reconstruct $r$ from its filtered versions by the pseudo-inverse
\[
\hat r=g\star\Big(\sum_n f_n^{\,r}\star\hat r_n\Big),\qquad
g\star\Big(\sum_n f_n^{\,r}\star f_n\Big)=\delta,
\]
where $f_n^{\,r}(x,y)=f_n(-x,-y)$ and $g$ is a precomputable, image-independent deconvolution filter. Then
$\hat l(x,y,t)=i(x,y,t)-\hat r(x,y)$.

\begin{algorithm}[H]
\SetAlgoLined
\KwIn{Log frames $i(\cdot,\cdot,t)$, $t=1..T$; derivative filters $\{f_n\}$ (usually $\partial_x,\partial_y$)}
\KwOut{Reflectance $\hat r$; illuminations $\hat l(\cdot,\cdot,t)$}
\ForEach{filter $f_n$}{$o_n(\cdot,\cdot,t)=i(\cdot,\cdot,t)\star f_n$; \ $\hat r_n=\operatorname{median}_t o_n$\;}
$\hat r\leftarrow g\star(\sum_n f_n^{\,r}\star\hat r_n)$ \tcp*{$g$ precomputed from the filters}
$\hat l(\cdot,\cdot,t)\leftarrow i(\cdot,\cdot,t)-\hat r$\;
\caption{Weiss intrinsic-image estimation}
\end{algorithm}

\textbf{Convergence (Claim 2).} If $p_\epsilon=P(|f_n\star l|<\epsilon)$ then
$P(|\hat r_n-r_n|<\epsilon)\ge\sum_{k=0}^{\lfloor T/2\rfloor}\binom{T}{k}(1-p_\epsilon)^{T-k}p_\epsilon^{\,k}$;
sparser lighting (larger $p_\epsilon$) converges faster --- $p_\epsilon=0.85$ gives $\ge0.93$ at $T=3$ and
$\ge0.97$ at $T=5$.

\Knobs
\begin{ptable}
Filters $\{f_n\}$ & Which gradients (usually $\partial_x,\partial_y$) & Two derivative filters are enough; $g$ depends only on them \\
\# frames $T$ & Convergence certainty & More frames / sparser light $\Rightarrow$ higher accuracy \\
$g$ & Reconstruction deconvolution & Precompute once; reuse on every sequence \\
\end{ptable}
\ShortExp Take logs so image $=$ paint $+$ light. The paint is constant across frames; the light's edges are
sparse (mostly zero). So at each pixel, the temporal \emph{median of the edge-filtered frames} throws away
the rare shadow-edge outliers and keeps the paint's gradients; reconstruct an image from those gradients (a
fixed deconvolution) to get reflectance, then subtract to get each frame's lighting. Crucially it is the
median of \emph{gradients}, not raw pixels --- a permanently-shadowed pixel produces no edge and cancels,
whereas a raw-pixel median (or min/mean) would bake the shadow into the paint.

% ##################################################################
\newpage
\section*{Course-wide Pattern}
\addcontentsline{toc}{section}{Course-wide Pattern}
% ##################################################################
\begin{notebox}
\textbf{Course-wide pattern.} Almost every method reduces to: \emph{define a cost, minimise it iteratively,
add a regulariser}. Least squares handles over-determined systems (motion models, LK/HS, super-resolution,
texture $A$); a smoothness/MAP prior resolves ambiguity (HS, super-resolution, colorization); dynamic
programming gives shortest paths (geodesic/Dijkstra, hyperlapse); predict/update or
propagation/randomisation drive the iterative solvers (Kalman, particle filter, PatchMatch, K-means); and
low-rank/SVD structure recurs (texture synthesis, fast KDE, intrinsic-image reconstruction).
\end{notebox}

% ##################################################################
\newpage
\appendix
\section*{Parameter Index and Glossary}
\addcontentsline{toc}{section}{Parameter Index and Glossary}
% ##################################################################
Quick reference: every parameter mentioned in this document, organised by algorithm. Use this to
quickly find what a parameter does or how changing it affects the method.

\subsection*{Affine Transformation (2\texttimes3)}
\textbf{$A$ (4 params)} --- Rotation/scale/shear

\textbf{$t=(t_x,t_y)$} --- Translation

\textbf{\# point pairs} --- Determinacy of the fit

\subsection*{Applications: Inpainting, Reshuffling, Object Removal}
\textbf{Mask $\Omega_m$} --- Region to synthesise

\textbf{Weights $W_P$} --- Preserve vs.\ remove ($\approx0$: erase; high: protect)

\textbf{Source set} --- Where fill comes from (multiple sources allowed)

\subsection*{BRIEF Descriptor}
\textbf{\# pixel pairs} --- Descriptor length (more: more distinctive, slower)

\textbf{Window size} --- Spatial support (larger: more context, blurrier)

\textbf{Pair pattern} --- Fixed sampling layout

\subsection*{Bayesian Recursive Tracking (derivation)}
\textbf{$p(x_t \mid x_{t-1})$} --- Motion / transition model

\textbf{$p(z_t \mid x_t)$} --- Measurement likelihood

\textbf{prior (recursion)} --- Carried-forward belief

\subsection*{Bidirectional-Similarity Re-targeting}
\textbf{$\alpha$ (0--1)} --- Completeness vs.\ coherence (higher: keep more content; lower: smoother)

\textbf{Patch size (5--7)} --- Context vs.\ sharpness

\textbf{Constraint weights $W_P$} --- Preserve / remove (high: protect a face/line; $\approx0$: erase)

\textbf{Pyramid step ($\sim1.1\times$)} --- Gradual resize for good initialisation

\subsection*{Dijkstra's Algorithm}
\textbf{Edge weights} --- Cost of crossing gradients (gradient magnitude)

\textbf{Neighbourhood $N_1$} --- 4-conn (image) or 6-conn (video with time)

\subsection*{Dynamic Texture Synthesis}
\textbf{$n$ (state dim)} --- Rank kept from SVD ($\ge95\%$ variance)

\textbf{State noise $v$} --- Trajectory exploration (prevents determinism)

\textbf{Image noise $w$} --- Natural look (hides low-rank seams)

\subsection*{Eulerian Video Magnification}
\textbf{$\alpha$ (amplification)} --- Motion gain, typically 5--50

\textbf{Band-pass band} --- Target frequency range (match your signal)

\textbf{Pyramid} --- Handles motions $>1$ pixel

\subsection*{Fast KDE via Low-Rank Kernel}
\textbf{$p$ (effective rank)} --- Approximation quality vs.\ speed

\textbf{Kernel choice} --- Pick low-rank kernels (fast-decaying singular values)

\subsection*{Fast Sweeping (Danielsson)}
\textbf{\# sweeps} --- Directions / turns covered (4 for simple paths; more for spirals)

\textbf{Neighbourhood layout} --- Enables parallel column-wise updates

\subsection*{Geodesic Distance Colorization}
\textbf{$b \in [1,6]$ (in $1/D^b$)} --- Transition sharpness (larger: sharper boundaries)

\textbf{Edge weight $|Y(p)-Y(q)|$} --- How edges act as walls

\textbf{Scribble placement} --- One per distinct FG/BG colour region (inside regions, not on edges)

\textbf{Video neighbourhood} --- 4 (image only) vs.\ 6 (+ time; lets colour flow through frames)

\subsection*{GMM Background Subtraction}
\textbf{$K$ (Gaussians)} --- Modes modelled, typically 3--5

\textbf{$\alpha$ (weight rate)} --- How fast new patterns become background ($\approx0.01$)

\textbf{$T$ (BG threshold)} --- Background fraction, 0.7--0.9

\textbf{Match distance ($2.5\sigma$)} --- Detection sensitivity to new objects

\subsection*{Harris Corner Detector}
\textbf{$\alpha$ (0.04--0.06)} --- Edge penalty in Harris response $R$

\textbf{Threshold $\theta$} --- Minimum cornerness (higher: fewer, stronger corners)

\textbf{Window $g$} --- Region summed into $H$ (larger: smoother, rotation-robust)

\subsection*{Homography (3\texttimes3)}
\textbf{$H$ (9 entries, 8 free)} --- Perspective warp parameters

\textbf{\# point pairs} --- Determinacy (need $\ge4$ pairs)

\textbf{Normalisation} --- Numerical stability (constrains overall scale)

\subsection*{Horn--Schunck}
\textbf{$\alpha$ (smoothness)} --- Data vs.\ prior; ML (small) vs.\ MAP (large)

\textbf{$K$ (iterations)} --- Solver convergence (more: closer to exact)

\subsection*{Hyperlapse}
\textbf{$\mu$ / $v$ (speedup)} --- Target frames skipped (larger: faster video)

\textbf{$\lambda_s$ (speedup weight)} --- Adherence to target rate (usually $\gg\lambda_a$)

\textbf{$\lambda_a$ (accel.\ weight)} --- Pace uniformity (larger: more constant intervals)

\textbf{window $w$} --- Max jump considered in DP (larger: more options)

\subsection*{Kalman Filter}
\textbf{$Q$ (process noise)} --- Trust in motion model (too small: over-confident; too big: jumpy)

\textbf{$R$ (measurement noise)} --- Trust in sensor ($R$ small: follow data; large: follow model)

\textbf{State vector} --- What is tracked (add velocity/accel for occlusion coasting)

\subsection*{K-means Clustering}
\textbf{$k$ (\# clusters)} --- Granularity (larger $k$: lower distortion but more noise)

\textbf{Initialisation} --- Which local minimum (k-means++ avoids empty clusters)

\textbf{Restarts} --- Robustness (keep the lowest-distortion run)

\subsection*{Layer Decomposition}
\textbf{$k$ (\# layers)} --- Number of motion groups (too small: lumps motions; too large: splits objects)

\textbf{Motion model} --- Affine (6) vs.\ homography (8)

\textbf{\# random windows} --- Sampling of motions (more: accurate, slower)

\subsection*{Linear Transformation (2\texttimes2)}
\textbf{$a,b,c,d$} --- Rotation/scale/shear jointly (4 free parameters)

\textbf{\# point pairs} --- Determinacy of the fit (need $\ge2$ pairs)

\subsection*{Lucas--Kanade}
\textbf{Window size $n$} --- \# equations per solve (larger: robust but mixes motions; smaller: worse aperture)

\textbf{Warp model} --- Translation (2) vs.\ affine (6)

\textbf{Pyramid levels} --- Max recoverable motion (more levels: handle bigger motion)

\subsection*{Matting (pure-foreground recovery)}
\textbf{Window size} --- Neighbour search for foreground/background (small windows usually suffice)

\textbf{New background $B$} --- Composite target (provided by user)

\subsection*{Opacity Map}
\textbf{$r \in (0,2]$} --- Distance weighting in opacity calculation

\textbf{Local KDE window} --- Boundary colour model (re-estimated from boundary pixels)

\subsection*{Opacity Map (Trimap Refinement)}
\textbf{$\rho$ (band radius)} --- Width of unsure band, 5--15 px

\textbf{$\sigma$ (KDE kernel)} --- Colour-model smoothness (larger: smoother distribution)

\textbf{Scribbles} --- Colour-region coverage (one per distinct FG/BG colour)

\subsection*{Particle Filter (Condensation)}
\textbf{$N$ (particles)} --- Approximation quality (more: better fidelity, costlier)

\textbf{Diffusion noise} --- Exploration; prevents sample impoverishment

\textbf{Output rule} --- Max vs.\ weighted mean (max: sharp but jittery; mean: stable)

\subsection*{PatchMatch}
\textbf{Patch size (5--7)} --- Context vs.\ sharpness

\textbf{$K$ iterations (4--6)} --- NN-field convergence

\textbf{Window $W$} --- Random-search radius (shrinks each pass)

\subsection*{Projection Matrix (pinhole camera)}
\textbf{$R,t$ (extrinsics)} --- Camera pose (6 params: 3 rotation, 3 translation)

\textbf{$f,k_u,k_v,u_0,v_0$ (intrinsics)} --- Lens + sensor mapping (5 params, calibrated once)

\textbf{Depth $Z$} --- Projects inversely to size ($1/Z$ shrinks distant objects)

\subsection*{RANSAC}
\textbf{$n$ (sample size)} --- Points to fit the model (minimal $n$ maximises all-inlier probability)

\textbf{$t$ (inlier threshold)} --- Inlier/outlier cutoff

\textbf{$K$ (iterations)} --- Success probability (set via the closed-form formula)

\subsection*{Re-targeting by Bidirectional Similarity}
\textbf{Initialisation $T^1$} --- Start from hierarchical pyramid, not black

\textbf{$\alpha$ (0--1)} --- Completeness vs.\ coherence (0.5 is a good start)

\subsection*{Super-Resolution (ML / MAP)}
\textbf{$\beta$ (step)} --- Gradient-descent rate

\textbf{$\lambda,\gamma$ (prior weights)} --- Smoothness vs.\ data in MAP

\textbf{\# frames} --- Recoverable detail (need $\gtrsim r^2$ with varied sub-pixel shifts)

\subsection*{Weiss Intrinsic Images}
\textbf{Filters $\{f_n\}$} --- Which gradients (usually $\partial_x,\partial_y$)

\textbf{\# frames $T$} --- Convergence certainty ($p_\epsilon=0.85$: $\ge0.93$ at $T=3$)

\textbf{$g$} --- Reconstruction deconvolution (precomputed from filters; reuse on every sequence)

\end{document}
